
\documentclass{article}
\usepackage{colortbl}
\usepackage{makecell}
\usepackage{multirow}
\usepackage{supertabular}

\begin{document}

\newcounter{utterance}

\centering \large Interaction Transcript for game `matchit\_ascii', experiment `different\_grid', episode 4 with mistral{-}large{-}3{-}azure.
\vspace{24pt}

{ \footnotesize  \setcounter{utterance}{1}
\setlength{\tabcolsep}{0pt}
\begin{supertabular}{c@{$\;$}|p{.15\linewidth}@{}p{.15\linewidth}p{.15\linewidth}p{.15\linewidth}p{.15\linewidth}p{.15\linewidth}}
    \# & $\;$A & \multicolumn{4}{c}{Game Master} & $\;\:$B\\
    \hline

    \theutterance \stepcounter{utterance}  
    & & \multicolumn{4}{p{0.6\linewidth}}{
        \cellcolor[rgb]{0.9,0.9,0.9}{
            \makecell[{{p{\linewidth}}}]{
                \texttt{\tiny{[P1$\langle$GM]}}
                \texttt{Participați la un joc colaborativ de ghicit.} \\
\texttt{Scopul este să aflați dacă această grilă și o altă imagine pe care doar eu o pot vedea sunt identice.} \\
\texttt{Vă rugăm să vă descrieți mai întâi grila. Apoi, eu voi furniza descrierea mea și ne putem pune reciproc întrebări despre grile pentru a afla dacă sunt identice. Acum începeți scurta dvs. descriere a grilei cu "DESCRIERE:", urmată de descriere. Nu adăugați nimic altceva. Aceasta este grila dumneavoastră:} \\
\texttt{X X X X X} \\
\texttt{▢ ▢ ▢ ▢ X} \\
\texttt{▢ ▢ X X X} \\
\texttt{▢ ▢ ▢ ▢ X} \\
\texttt{▢ ▢ ▢ ▢ X} \\
            }
        }
    }
    & \\ \\

    \theutterance \stepcounter{utterance}  
    & \multicolumn{4}{p{0.6\linewidth}}{
        \cellcolor[rgb]{0.9,0.9,0.9}{
            \makecell[{{p{\linewidth}}}]{
                \texttt{\tiny{[P1$\rangle$GM]}}
                \texttt{DESCRIERE: Grila este formată din 5 rânduri și 5 coloane. Simbolurile folosite sunt "X" și "▢". Primul rând conține doar "X" (5 simboluri). Al doilea rând are 4 "▢" urmate de un "X". Al treilea rând are 2 "▢", apoi 3 "X". Al patrulea și al cincilea rând au fiecare 4 "▢" urmate de un "X".} \\
            }
        }
    }
    & & \\ \\

    \theutterance \stepcounter{utterance}  
    & & & \multicolumn{2}{p{0.3\linewidth}}{
        \cellcolor[rgb]{0.9,0.9,0.9}{
            \makecell[{{p{\linewidth}}}]{
                \texttt{\tiny{[GM$|$GM]}}
                \texttt{continue} \\
            }
        }
    }
    & & \\ \\

    \theutterance \stepcounter{utterance}  
    & & \multicolumn{4}{p{0.6\linewidth}}{
        \cellcolor[rgb]{0.9,0.9,0.9}{
            \makecell[{{p{\linewidth}}}]{
                \texttt{\tiny{[GM$\rangle$P2]}}
                \texttt{Participați la un joc colaborativ de ghicit.} \\
\texttt{Scopul este să aflați dacă această grilă și o altă imagine pe care doar eu o pot vedea sunt identice.} \\
\texttt{Vă rugăm să vă descrieți mai întâi grila. Apoi, eu voi furniza descrierea mea și ne putem pune reciproc întrebări despre grile pentru a afla dacă sunt identice. Acum începeți scurta dvs. descriere a grilei cu "DESCRIERE:", urmată de descriere. Nu adăugați nimic altceva. Aceasta este grila dumneavoastră:} \\
\texttt{▢ ▢ ▢ ▢ ▢} \\
\texttt{▢ X X X ▢} \\
\texttt{▢ X X X ▢} \\
\texttt{▢ X X X ▢} \\
\texttt{▢ ▢ ▢ ▢ ▢} \\
            }
        }
    }
    & \\ \\

    \theutterance \stepcounter{utterance}  
    & & & \multicolumn{4}{p{0.6\linewidth}}{
        \cellcolor[rgb]{0.9,0.9,0.9}{
            \makecell[{{p{\linewidth}}}]{
                \texttt{\tiny{[GM$\langle$P2]}}
                \texttt{DESCRIERE: Grila este formată din 5 rânduri și 5 coloane, alcătuind un pătrat. Toate celulele de pe margini (primul și ultimul rând, prima și ultima coloană) sunt goale (▢). În centrul grilei, există un bloc compact de 3x3 celule marcate cu "X".} \\
            }
        }
    }
     \\ \\

    \theutterance \stepcounter{utterance}  
    & & & \multicolumn{2}{p{0.3\linewidth}}{
        \cellcolor[rgb]{0.9,0.9,0.9}{
            \makecell[{{p{\linewidth}}}]{
                \texttt{\tiny{[GM$|$GM]}}
                \texttt{continue} \\
            }
        }
    }
    & & \\ \\

    \theutterance \stepcounter{utterance}  
    & & \multicolumn{4}{p{0.6\linewidth}}{
        \cellcolor[rgb]{0.9,0.9,0.9}{
            \makecell[{{p{\linewidth}}}]{
                \texttt{\tiny{[GM$\rangle$P2]}}
                \texttt{Aceasta este DESCRIEREA meaDESCRIERE: Grila este formată din 5 rânduri și 5 coloane. Simbolurile folosite sunt "X" și "▢". Primul rând conține doar "X" (5 simboluri). Al doilea rând are 4 "▢" urmate de un "X". Al treilea rând are 2 "▢", apoi 3 "X". Al patrulea și al cincilea rând au fiecare 4 "▢" urmate de un "X".} \\
\texttt{Acum pune o întrebare pentru a afla noi aspecte ale grilei mele care ar putea fi diferite de grila ta. Începe cu "ÎNTREBARE:" și nu adăuga nimic altceva.} \\
            }
        }
    }
    & \\ \\

    \theutterance \stepcounter{utterance}  
    & & & \multicolumn{4}{p{0.6\linewidth}}{
        \cellcolor[rgb]{0.9,0.9,0.9}{
            \makecell[{{p{\linewidth}}}]{
                \texttt{\tiny{[GM$\langle$P2]}}
                \texttt{ÎNTREBARE: Există vreun "X" în celulele de pe marginea exterioară a grilei (primul/ultimul rând sau prima/ultima coloană), în afara centrului 3x3?} \\
            }
        }
    }
     \\ \\

    \theutterance \stepcounter{utterance}  
    & & & \multicolumn{2}{p{0.3\linewidth}}{
        \cellcolor[rgb]{0.9,0.9,0.9}{
            \makecell[{{p{\linewidth}}}]{
                \texttt{\tiny{[GM$|$GM]}}
                \texttt{continue} \\
            }
        }
    }
    & & \\ \\

    \theutterance \stepcounter{utterance}  
    & & \multicolumn{4}{p{0.6\linewidth}}{
        \cellcolor[rgb]{0.9,0.9,0.9}{
            \makecell[{{p{\linewidth}}}]{
                \texttt{\tiny{[P1$\langle$GM]}}
                \texttt{Aceasta este DESCRIEREA meaDESCRIERE: Grila este formată din 5 rânduri și 5 coloane, alcătuind un pătrat. Toate celulele de pe margini (primul și ultimul rând, prima și ultima coloană) sunt goale (▢). În centrul grilei, există un bloc compact de 3x3 celule marcate cu "X".} \\
\texttt{ÎNTREBARE: Există vreun "X" în celulele de pe marginea exterioară a grilei (primul/ultimul rând sau prima/ultima coloană), în afara centrului 3x3?Începeți răspunsul cu "RĂSPUNS:" și nu adăugați nimic altceva.} \\
            }
        }
    }
    & \\ \\

    \theutterance \stepcounter{utterance}  
    & \multicolumn{4}{p{0.6\linewidth}}{
        \cellcolor[rgb]{0.9,0.9,0.9}{
            \makecell[{{p{\linewidth}}}]{
                \texttt{\tiny{[P1$\rangle$GM]}}
                \texttt{RĂSPUNS: Da, există "X" în primul rând (toate celulele) și în ultima coloană (rândurile 2, 3, 4 și 5).} \\
            }
        }
    }
    & & \\ \\

    \theutterance \stepcounter{utterance}  
    & & & \multicolumn{2}{p{0.3\linewidth}}{
        \cellcolor[rgb]{0.9,0.9,0.9}{
            \makecell[{{p{\linewidth}}}]{
                \texttt{\tiny{[GM$|$GM]}}
                \texttt{continue} \\
            }
        }
    }
    & & \\ \\

    \theutterance \stepcounter{utterance}  
    & & \multicolumn{4}{p{0.6\linewidth}}{
        \cellcolor[rgb]{0.9,0.9,0.9}{
            \makecell[{{p{\linewidth}}}]{
                \texttt{\tiny{[P1$\langle$GM]}}
                \texttt{Acum pune o întrebare pentru a afla noi aspecte ale grilei mele care ar putea fi diferite de grila ta. Începe cu "ÎNTREBARE:" și nu adăuga nimic altceva.} \\
            }
        }
    }
    & \\ \\

    \theutterance \stepcounter{utterance}  
    & \multicolumn{4}{p{0.6\linewidth}}{
        \cellcolor[rgb]{0.9,0.9,0.9}{
            \makecell[{{p{\linewidth}}}]{
                \texttt{\tiny{[P1$\rangle$GM]}}
                \texttt{ÎNTREBARE: Există vreo celulă marcată cu "X" în afara blocului central 3x3 (excluzând marginile, dacă este cazul) sau toate "X"{-}urile sunt concentrate strict în centrul grilei?} \\
            }
        }
    }
    & & \\ \\

    \theutterance \stepcounter{utterance}  
    & & & \multicolumn{2}{p{0.3\linewidth}}{
        \cellcolor[rgb]{0.9,0.9,0.9}{
            \makecell[{{p{\linewidth}}}]{
                \texttt{\tiny{[GM$|$GM]}}
                \texttt{continue} \\
            }
        }
    }
    & & \\ \\

    \theutterance \stepcounter{utterance}  
    & & \multicolumn{4}{p{0.6\linewidth}}{
        \cellcolor[rgb]{0.9,0.9,0.9}{
            \makecell[{{p{\linewidth}}}]{
                \texttt{\tiny{[GM$\rangle$P2]}}
                \texttt{RĂSPUNS: Da, există "X" în primul rând (toate celulele) și în ultima coloană (rândurile 2, 3, 4 și 5).} \\
\texttt{ÎNTREBARE: Există vreo celulă marcată cu "X" în afara blocului central 3x3 (excluzând marginile, dacă este cazul) sau toate "X"{-}urile sunt concentrate strict în centrul grilei?Începeți răspunsul cu "RĂSPUNS:" și nu adăugați nimic altceva.} \\
            }
        }
    }
    & \\ \\

    \theutterance \stepcounter{utterance}  
    & & & \multicolumn{4}{p{0.6\linewidth}}{
        \cellcolor[rgb]{0.9,0.9,0.9}{
            \makecell[{{p{\linewidth}}}]{
                \texttt{\tiny{[GM$\langle$P2]}}
                \texttt{RĂSPUNS: Toate "X"{-}urile din grila mea sunt concentrate strict în blocul central 3x3, fără alte "X" în interiorul grilei în afara acestuia.} \\
            }
        }
    }
     \\ \\

    \theutterance \stepcounter{utterance}  
    & & & \multicolumn{2}{p{0.3\linewidth}}{
        \cellcolor[rgb]{0.9,0.9,0.9}{
            \makecell[{{p{\linewidth}}}]{
                \texttt{\tiny{[GM$|$GM]}}
                \texttt{continue} \\
            }
        }
    }
    & & \\ \\

    \theutterance \stepcounter{utterance}  
    & & \multicolumn{4}{p{0.6\linewidth}}{
        \cellcolor[rgb]{0.9,0.9,0.9}{
            \makecell[{{p{\linewidth}}}]{
                \texttt{\tiny{[GM$\rangle$P2]}}
                \texttt{Acum pune o întrebare pentru a afla noi aspecte ale grilei mele care ar putea fi diferite de grila ta. Începe cu "ÎNTREBARE:" și nu adăuga nimic altceva.} \\
            }
        }
    }
    & \\ \\

    \theutterance \stepcounter{utterance}  
    & & & \multicolumn{4}{p{0.6\linewidth}}{
        \cellcolor[rgb]{0.9,0.9,0.9}{
            \makecell[{{p{\linewidth}}}]{
                \texttt{\tiny{[GM$\langle$P2]}}
                \texttt{ÎNTREBARE: Care este poziția exactă a singurului "X" din rândurile 2, 4 și 5? (De exemplu: coloana 1, 2, 3, 4 sau 5?)} \\
            }
        }
    }
     \\ \\

    \theutterance \stepcounter{utterance}  
    & & & \multicolumn{2}{p{0.3\linewidth}}{
        \cellcolor[rgb]{0.9,0.9,0.9}{
            \makecell[{{p{\linewidth}}}]{
                \texttt{\tiny{[GM$|$GM]}}
                \texttt{continue} \\
            }
        }
    }
    & & \\ \\

    \theutterance \stepcounter{utterance}  
    & & \multicolumn{4}{p{0.6\linewidth}}{
        \cellcolor[rgb]{0.9,0.9,0.9}{
            \makecell[{{p{\linewidth}}}]{
                \texttt{\tiny{[P1$\langle$GM]}}
                \texttt{RĂSPUNS: Toate "X"{-}urile din grila mea sunt concentrate strict în blocul central 3x3, fără alte "X" în interiorul grilei în afara acestuia.} \\
\texttt{ÎNTREBARE: Care este poziția exactă a singurului "X" din rândurile 2, 4 și 5? (De exemplu: coloana 1, 2, 3, 4 sau 5?)Începeți răspunsul cu "RĂSPUNS:" și nu adăugați nimic altceva.} \\
            }
        }
    }
    & \\ \\

    \theutterance \stepcounter{utterance}  
    & \multicolumn{4}{p{0.6\linewidth}}{
        \cellcolor[rgb]{0.9,0.9,0.9}{
            \makecell[{{p{\linewidth}}}]{
                \texttt{\tiny{[P1$\rangle$GM]}}
                \texttt{RĂSPUNS: În rândurile 2, 4 și 5, singurul "X" este poziționat în **coloana 5** (ultima coloană).} \\
            }
        }
    }
    & & \\ \\

    \theutterance \stepcounter{utterance}  
    & & & \multicolumn{2}{p{0.3\linewidth}}{
        \cellcolor[rgb]{0.9,0.9,0.9}{
            \makecell[{{p{\linewidth}}}]{
                \texttt{\tiny{[GM$|$GM]}}
                \texttt{continue} \\
            }
        }
    }
    & & \\ \\

    \theutterance \stepcounter{utterance}  
    & & \multicolumn{4}{p{0.6\linewidth}}{
        \cellcolor[rgb]{0.9,0.9,0.9}{
            \makecell[{{p{\linewidth}}}]{
                \texttt{\tiny{[P1$\langle$GM]}}
                \texttt{Acum pune o întrebare pentru a afla noi aspecte ale grilei mele care ar putea fi diferite de grila ta. Începe cu "ÎNTREBARE:" și nu adăuga nimic altceva.} \\
            }
        }
    }
    & \\ \\

    \theutterance \stepcounter{utterance}  
    & \multicolumn{4}{p{0.6\linewidth}}{
        \cellcolor[rgb]{0.9,0.9,0.9}{
            \makecell[{{p{\linewidth}}}]{
                \texttt{\tiny{[P1$\rangle$GM]}}
                \texttt{ÎNTREBARE: Blocul central 3x3 din grila ta este complet umplut cu "X" (fără niciun "▢" în interiorul lui), sau există celule goale în interiorul acestuia?} \\
            }
        }
    }
    & & \\ \\

    \theutterance \stepcounter{utterance}  
    & & & \multicolumn{2}{p{0.3\linewidth}}{
        \cellcolor[rgb]{0.9,0.9,0.9}{
            \makecell[{{p{\linewidth}}}]{
                \texttt{\tiny{[GM$|$GM]}}
                \texttt{continue} \\
            }
        }
    }
    & & \\ \\

    \theutterance \stepcounter{utterance}  
    & & \multicolumn{4}{p{0.6\linewidth}}{
        \cellcolor[rgb]{0.9,0.9,0.9}{
            \makecell[{{p{\linewidth}}}]{
                \texttt{\tiny{[GM$\rangle$P2]}}
                \texttt{RĂSPUNS: În rândurile 2, 4 și 5, singurul "X" este poziționat în **coloana 5** (ultima coloană).} \\
\texttt{ÎNTREBARE: Blocul central 3x3 din grila ta este complet umplut cu "X" (fără niciun "▢" în interiorul lui), sau există celule goale în interiorul acestuia?Începeți răspunsul cu "RĂSPUNS:" și nu adăugați nimic altceva.} \\
            }
        }
    }
    & \\ \\

    \theutterance \stepcounter{utterance}  
    & & & \multicolumn{4}{p{0.6\linewidth}}{
        \cellcolor[rgb]{0.9,0.9,0.9}{
            \makecell[{{p{\linewidth}}}]{
                \texttt{\tiny{[GM$\langle$P2]}}
                \texttt{RĂSPUNS: Blocul central 3x3 din grila mea este complet umplut cu "X" (fără niciun "▢" în interiorul lui).} \\
            }
        }
    }
     \\ \\

    \theutterance \stepcounter{utterance}  
    & & & \multicolumn{2}{p{0.3\linewidth}}{
        \cellcolor[rgb]{0.9,0.9,0.9}{
            \makecell[{{p{\linewidth}}}]{
                \texttt{\tiny{[GM$|$GM]}}
                \texttt{continue} \\
            }
        }
    }
    & & \\ \\

    \theutterance \stepcounter{utterance}  
    & & \multicolumn{4}{p{0.6\linewidth}}{
        \cellcolor[rgb]{0.9,0.9,0.9}{
            \makecell[{{p{\linewidth}}}]{
                \texttt{\tiny{[GM$\rangle$P2]}}
                \texttt{Acum pune o întrebare pentru a afla noi aspecte ale grilei mele care ar putea fi diferite de grila ta. Începe cu "ÎNTREBARE:" și nu adăuga nimic altceva.} \\
            }
        }
    }
    & \\ \\

    \theutterance \stepcounter{utterance}  
    & & & \multicolumn{4}{p{0.6\linewidth}}{
        \cellcolor[rgb]{0.9,0.9,0.9}{
            \makecell[{{p{\linewidth}}}]{
                \texttt{\tiny{[GM$\langle$P2]}}
                \texttt{ÎNTREBARE: În rândul 3 al grilei mele, cele două "▢" sunt plasate consecutiv (una lângă alta) sau sunt separate de "X"{-}uri? (De exemplu: ▢▢XXX sau ▢X▢XX etc.)} \\
            }
        }
    }
     \\ \\

    \theutterance \stepcounter{utterance}  
    & & & \multicolumn{2}{p{0.3\linewidth}}{
        \cellcolor[rgb]{0.9,0.9,0.9}{
            \makecell[{{p{\linewidth}}}]{
                \texttt{\tiny{[GM$|$GM]}}
                \texttt{continue} \\
            }
        }
    }
    & & \\ \\

    \theutterance \stepcounter{utterance}  
    & & \multicolumn{4}{p{0.6\linewidth}}{
        \cellcolor[rgb]{0.9,0.9,0.9}{
            \makecell[{{p{\linewidth}}}]{
                \texttt{\tiny{[P1$\langle$GM]}}
                \texttt{RĂSPUNS: Blocul central 3x3 din grila mea este complet umplut cu "X" (fără niciun "▢" în interiorul lui).} \\
\texttt{ÎNTREBARE: În rândul 3 al grilei mele, cele două "▢" sunt plasate consecutiv (una lângă alta) sau sunt separate de "X"{-}uri? (De exemplu: ▢▢XXX sau ▢X▢XX etc.)Începeți răspunsul cu "RĂSPUNS:" și nu adăugați nimic altceva.} \\
            }
        }
    }
    & \\ \\

    \theutterance \stepcounter{utterance}  
    & \multicolumn{4}{p{0.6\linewidth}}{
        \cellcolor[rgb]{0.9,0.9,0.9}{
            \makecell[{{p{\linewidth}}}]{
                \texttt{\tiny{[P1$\rangle$GM]}}
                \texttt{RĂSPUNS: În rândul 3 al grilei mele, cele două "▢" sunt plasate consecutiv la început (▢▢XXX).} \\
            }
        }
    }
    & & \\ \\

    \theutterance \stepcounter{utterance}  
    & & & \multicolumn{2}{p{0.3\linewidth}}{
        \cellcolor[rgb]{0.9,0.9,0.9}{
            \makecell[{{p{\linewidth}}}]{
                \texttt{\tiny{[GM$|$GM]}}
                \texttt{continue} \\
            }
        }
    }
    & & \\ \\

    \theutterance \stepcounter{utterance}  
    & & \multicolumn{4}{p{0.6\linewidth}}{
        \cellcolor[rgb]{0.9,0.9,0.9}{
            \makecell[{{p{\linewidth}}}]{
                \texttt{\tiny{[P1$\langle$GM]}}
                \texttt{Acum pune o întrebare pentru a afla noi aspecte ale grilei mele care ar putea fi diferite de grila ta. Începe cu "ÎNTREBARE:" și nu adăuga nimic altceva.} \\
            }
        }
    }
    & \\ \\

    \theutterance \stepcounter{utterance}  
    & \multicolumn{4}{p{0.6\linewidth}}{
        \cellcolor[rgb]{0.9,0.9,0.9}{
            \makecell[{{p{\linewidth}}}]{
                \texttt{\tiny{[P1$\rangle$GM]}}
                \texttt{ÎNTREBARE: În grila ta, există vreun rând (în afara blocului central 3x3) care să conțină mai mult de un "X"? Dacă da, care este acesta?} \\
            }
        }
    }
    & & \\ \\

    \theutterance \stepcounter{utterance}  
    & & & \multicolumn{2}{p{0.3\linewidth}}{
        \cellcolor[rgb]{0.9,0.9,0.9}{
            \makecell[{{p{\linewidth}}}]{
                \texttt{\tiny{[GM$|$GM]}}
                \texttt{continue} \\
            }
        }
    }
    & & \\ \\

    \theutterance \stepcounter{utterance}  
    & & \multicolumn{4}{p{0.6\linewidth}}{
        \cellcolor[rgb]{0.9,0.9,0.9}{
            \makecell[{{p{\linewidth}}}]{
                \texttt{\tiny{[GM$\rangle$P2]}}
                \texttt{RĂSPUNS: În rândul 3 al grilei mele, cele două "▢" sunt plasate consecutiv la început (▢▢XXX).} \\
\texttt{ÎNTREBARE: În grila ta, există vreun rând (în afara blocului central 3x3) care să conțină mai mult de un "X"? Dacă da, care este acesta?Începeți răspunsul cu "RĂSPUNS:" și nu adăugați nimic altceva.} \\
            }
        }
    }
    & \\ \\

    \theutterance \stepcounter{utterance}  
    & & & \multicolumn{4}{p{0.6\linewidth}}{
        \cellcolor[rgb]{0.9,0.9,0.9}{
            \makecell[{{p{\linewidth}}}]{
                \texttt{\tiny{[GM$\langle$P2]}}
                \texttt{RĂSPUNS: Nu, în grila mea **niciun rând** (inclusiv marginile) nu conține mai mult de un "X" în afara blocului central 3x3. Singurele "X"{-}uri din afara centrului sunt cele de pe margini (primul rând complet și câte un "X" în coloana 5 pentru rândurile 2, 4 și 5).} \\
            }
        }
    }
     \\ \\

    \theutterance \stepcounter{utterance}  
    & & & \multicolumn{2}{p{0.3\linewidth}}{
        \cellcolor[rgb]{0.9,0.9,0.9}{
            \makecell[{{p{\linewidth}}}]{
                \texttt{\tiny{[GM$|$GM]}}
                \texttt{continue} \\
            }
        }
    }
    & & \\ \\

    \theutterance \stepcounter{utterance}  
    & & \multicolumn{4}{p{0.6\linewidth}}{
        \cellcolor[rgb]{0.9,0.9,0.9}{
            \makecell[{{p{\linewidth}}}]{
                \texttt{\tiny{[GM$\rangle$P2]}}
                \texttt{Acum ia o decizie. Ce crezi: grila ta și cealaltă grilă descrisă sunt la fel? Scrie "DECIZIE: aceeași grilă" dacă crezi că sunt aceeași grilă sau "DECIZIE: grile diferite" dacă crezi că sunt grile diferite. Nu adăuga nimic altceva.} \\
            }
        }
    }
    & \\ \\

    \theutterance \stepcounter{utterance}  
    & & & \multicolumn{4}{p{0.6\linewidth}}{
        \cellcolor[rgb]{0.9,0.9,0.9}{
            \makecell[{{p{\linewidth}}}]{
                \texttt{\tiny{[GM$\langle$P2]}}
                \texttt{DECIZIE: grile diferite} \\
            }
        }
    }
     \\ \\

    \theutterance \stepcounter{utterance}  
    & & & \multicolumn{2}{p{0.3\linewidth}}{
        \cellcolor[rgb]{0.9,0.9,0.9}{
            \makecell[{{p{\linewidth}}}]{
                \texttt{\tiny{[GM$|$GM]}}
                \texttt{continue} \\
            }
        }
    }
    & & \\ \\

    \theutterance \stepcounter{utterance}  
    & & & \multicolumn{2}{p{0.3\linewidth}}{
        \cellcolor[rgb]{0.9,0.9,0.9}{
            \makecell[{{p{\linewidth}}}]{
                \texttt{\tiny{[GM$|$GM]}}
                \texttt{success} \\
            }
        }
    }
    & & \\ \\

    \theutterance \stepcounter{utterance}  
    & & \multicolumn{4}{p{0.6\linewidth}}{
        \cellcolor[rgb]{0.9,0.9,0.9}{
            \makecell[{{p{\linewidth}}}]{
                \texttt{\tiny{[P1$\langle$GM]}}
                \texttt{RĂSPUNS: Nu, în grila mea **niciun rând** (inclusiv marginile) nu conține mai mult de un "X" în afara blocului central 3x3. Singurele "X"{-}uri din afara centrului sunt cele de pe margini (primul rând complet și câte un "X" în coloana 5 pentru rândurile 2, 4 și 5).} \\
\texttt{Acum ia o decizie. Ce crezi: grila ta și cealaltă grilă descrisă sunt la fel? Scrie "DECIZIE: aceeași grilă" dacă crezi că sunt aceeași grilă sau "DECIZIE: grile diferite" dacă crezi că sunt grile diferite. Nu adăuga nimic altceva.} \\
            }
        }
    }
    & \\ \\

    \theutterance \stepcounter{utterance}  
    & \multicolumn{4}{p{0.6\linewidth}}{
        \cellcolor[rgb]{0.9,0.9,0.9}{
            \makecell[{{p{\linewidth}}}]{
                \texttt{\tiny{[P1$\rangle$GM]}}
                \texttt{DECIZIE: grile diferite} \\
            }
        }
    }
    & & \\ \\

    \theutterance \stepcounter{utterance}  
    & & & \multicolumn{2}{p{0.3\linewidth}}{
        \cellcolor[rgb]{0.9,0.9,0.9}{
            \makecell[{{p{\linewidth}}}]{
                \texttt{\tiny{[GM$|$GM]}}
                \texttt{continue} \\
            }
        }
    }
    & & \\ \\

    \theutterance \stepcounter{utterance}  
    & & & \multicolumn{2}{p{0.3\linewidth}}{
        \cellcolor[rgb]{0.9,0.9,0.9}{
            \makecell[{{p{\linewidth}}}]{
                \texttt{\tiny{[GM$|$GM]}}
                \texttt{success} \\
            }
        }
    }
    & & \\ \\

\end{supertabular}
}

\end{document}
